\begin{mistake}{2026-07-19}{03}

\begin{problemcontent}
设 \(f(x)\) 在 \((-\infty, +\infty)\) 内连续，且满足 \(f(x+T) = f(x)\) (\(T>0\))，\(f(-x) = f(x)\)。
证明：\(\int_0^{nT} x f(x) dx = \frac{n^2 T}{2} \int_0^T f(x) dx\) (\(n\) 为正整数)。
\end{problemcontent}

\begin{solutioncontent}
设 \(I = \int_0^{nT} x f(x) dx\)。
令 \(x = nT - t\)，则 \(dx = -dt\)，积分上下限由 \([0, nT]\) 变为 \([nT, 0]\)。
代入原积分并吸收负号调换上下限：
\[I = \int_0^{nT} (nT - t) f(nT - t) dt\]
由 \(f(x)\) 以 \(T\) 为周期及 \(f(-x) = f(x)\) 的偶函数性质，有：
\[f(nT - t) = f(-t) = f(t)\]
将其代入积分并展开：
\[I = \int_0^{nT} (nT - t) f(t) dt = nT \int_0^{nT} f(t) dt - \int_0^{nT} t f(t) dt\]
注意到等式最右侧第二项即为原积分 \(I\)（积分变量符号不影响积分值），故移项得：
\[2I = nT \int_0^{nT} f(t) dt \implies I = \frac{nT}{2} \int_0^{nT} f(t) dt\]
又根据周期函数在整周期上的积分性质，有 \(\int_0^{nT} f(t) dt = n \int_0^T f(t) dt\)。
代入上式即得：
\[I = \frac{nT}{2} \left( n \int_0^T f(t) dt \right) = \frac{n^2 T}{2} \int_0^T f(x) dx\]
证明完毕。
\end{solutioncontent}

\mistakeknowledge{
\begin{itemize}
  \item \textbf{核心公式/定理}：区间翻转（再现）代换 \(x = a+b-t\)；周期函数积分性质 \(\int_0^{nT} f(x) dx = n \int_0^T f(x) dx\)。
  \item \textbf{易错点}：化简 \(f(nT-t)\) 时容易受阻，需明确：周期性负责消除 \(nT\)，奇偶性负责消除内部负号。
  \item \textbf{关联知识}：带有 \(x f(x)\) 结构的定积分，优先尝试区间翻转代换构造方程消去 \(x\)，比分段积分求和更高效。
\end{itemize}}

\end{mistake}
